\subsection{解析全息反演、二次提升账本与时间圆维下界}\label{subsec:conclusion-analytic-holography-joukowsky-ledger-time-dimension}
为将解析全息几何刚性、Joukowsky--Godel 输运的代数分歧、window-\(m\) 压缩率与时间视窗圆维下界统一在同一可审计链条中，本小节给出一组可独立引用的结论及证明。

\leanverified{paper\_conclusion\_analytic\_holography\_inversion\_rigidity}
\begin{theorem}[解析全息反演刚性与二阶径向矩的唯一读出]\label{thm:conclusion-analytic-holography-inversion-rigidity}
设 \(\Gamma\subset\CC\) 为实解析 Jordan 曲线，\(\CC\setminus\Gamma=\Omega^0\sqcup\Omega^\infty\)（\(\infty\in\Omega^\infty\)）。
设 \(\mu\) 为支撑于 \(\Gamma\) 的概率测度，其 Cauchy 变换
\[
G_\mu(z):=\int_\Gamma \frac{1}{z-w}\,\dd\mu(w)
\]
满足
\[
G_\mu(z)=0\quad(z\in\Omega^0),\qquad
G_\mu(z)=\frac{1}{\sqrt{z^2-4}}\quad(z\in\Omega^\infty),
\]
其中分支取为 \(\sqrt{z^2-4}\sim z\)（\(z\to\infty\)）。
则存在唯一 \(r\ge 1\) 使得
\[
\Gamma=E_r:=\left\{r e^{i\theta}+r^{-1}e^{-i\theta}:\ \theta\in[0,2\pi)\right\},
\]
且
\[
\mu=\mu_r:=(J_r)_\# m,\qquad J_r(\eta):=r\eta+r^{-1}\eta^{-1},
\]
其中 \(m\) 为单位圆 \(\TT\) 上的归一化 Haar 测度。
此外，定义
\[
R_2(\mu):=\int_\Gamma |w|^2\,\dd\mu(w),
\]
则
\[
R_2(\mu)=r^2+r^{-2},
\]
并且参数可由
\[
r^2=\frac{R_2(\mu)+\sqrt{R_2(\mu)^2-4}}{2},
\qquad
r=\left(\frac{R_2(\mu)+\sqrt{R_2(\mu)^2-4}}{2}\right)^{1/2}
\]
唯一恢复。
\end{theorem}
\begin{proof}
在 \(\Omega^\infty\) 上定义
\[
\Phi(z):=\frac{z+\sqrt{z^2-4}}{2},
\]
则直接计算得
\[
\frac{\dd}{\dd z}\log\Phi(z)=\frac{1}{\sqrt{z^2-4}}=G_\mu(z)\qquad(z\in\Omega^\infty).
\]
取原函数 \(F\) 使 \(F'(z)=G_\mu(z)\)。
由 \(G_\mu\equiv 0\) 于 \(\Omega^0\)，\(F\) 在 \(\Omega^0\) 为常数；
在 \(\Omega^\infty\) 则有 \(F(z)=\log\Phi(z)+C\)。
令 \(U:=\Re F\)，则 \(U\) 在 \(\Omega^0\) 常值，且在 \(\Omega^\infty\) 满足
\[
U(z)=\log|\Phi(z)|+\Re C.
\]
由边界连续性，\(\Gamma=\partial\Omega^0\) 必位于某条等势线 \(|\Phi|=\rho\)（\(\rho>0\)）上。
记 \(\zeta=\Phi(z)\)，利用恒等式 \(z=\zeta+\zeta^{-1}\)，当 \(|\zeta|=\rho\) 时可写
\[
z=\rho\eta+\rho^{-1}\eta^{-1},\qquad \eta\in\TT.
\]
故 \(\Gamma=E_\rho\)，又 \(E_\rho=E_{\rho^{-1}}\)，取 \(r:=\max\{\rho,\rho^{-1}\}\ge 1\) 得唯一参数化。

设 \(\mu_r=(J_r)_\#m\)。
由 \(J_r\) 的拉回积分与留数计算，得其 Cauchy 变换满足
\[
G_{\mu_r}(z)=0\quad(z\in\Omega^0),\qquad
G_{\mu_r}(z)=\frac{1}{\sqrt{z^2-4}}\quad(z\in\Omega^\infty).
\]
于是 \(G_{\mu-\mu_r}\equiv 0\) 于 \(\CC\setminus\Gamma\) 两连通分支，按 Sokhotski--Plemelj 跳跃公式可知 \(\mu-\mu_r=0\)，即 \(\mu=\mu_r\)。

最后，按参数化 \(w(\theta)=r e^{i\theta}+r^{-1}e^{-i\theta}\) 计算
\[
R_2(\mu)=\frac{1}{2\pi}\int_0^{2\pi}|w(\theta)|^2\,\dd\theta
=\frac{1}{2\pi}\int_0^{2\pi}\!\left(r^2+r^{-2}+e^{2i\theta}+e^{-2i\theta}\right)\dd\theta
=r^2+r^{-2}.
\]
令 \(s=r^2\ge 1\)，由 \(s+s^{-1}=R_2(\mu)\) 得
\[
s^2-R_2(\mu)s+1=0,
\]
故
\[
s=\frac{R_2(\mu)+\sqrt{R_2(\mu)^2-4}}{2},
\]
从而得到所述恢复公式与唯一性。证毕。
\end{proof}

\leanverified{paper\_conclusion\_jg\_slit\_uniformization\_orthogonal\_coordinates}
\begin{theorem}[裂平面均匀化与共焦椭圆--双曲正交坐标]\label{thm:conclusion-jg-slit-uniformization-orthogonal-coordinates}
令
\[
\Omega_{\mathrm{slit}}:=\CC\setminus[-2,2],
\]
并在 \(\Omega_{\mathrm{slit}}\) 上取平方根分支 \(\sqrt{w^2-4}\sim w\)（\(w\to\infty\)）。
定义均匀化映射
\[
z(w):=\frac{w+\sqrt{w^2-4}}{2},
\qquad |z(w)|>1.
\]
则对任意 \(w\in\Omega_{\mathrm{slit}}\)，存在唯一 \((u,\theta)\in(0,\infty)\times(-\pi,\pi]\) 使得
\[
z(w)=e^{u+i\theta},
\qquad
w=e^{u+i\theta}+e^{-u-i\theta}.
\]
写 \(w=x+iy\)，则
\[
x=2\cosh u\cos\theta,\qquad
y=2\sinh u\sin\theta.
\]
并且：
\begin{enumerate}
\normalfont
  \item \(u=\mathrm{const}\) 的坐标曲线为以 \(\pm2\) 为焦点的共焦椭圆族；
  \item \(\theta=\mathrm{const}\) 的坐标曲线为以 \(\pm2\) 为焦点的共焦双曲线族（当 \(\sin\theta\cos\theta\neq 0\) 时为非退化双曲线）；
  \item \((u,\theta)\) 为欧氏度量下的正交坐标，并满足共形因子公式
  \[
  \dd s^2
  =4\bigl(\cosh^2u-\cos^2\theta\bigr)\,(\dd u^2+\dd\theta^2).
  \]
\end{enumerate}
\end{theorem}
\begin{proof}
由二次方程 \(z^2-wz+1=0\) 的两根为 \(z\) 与 \(z^{-1}\)，分支选取保证 \(z(w)\) 为满足 \(|z|>1\) 的根，从而唯一表示为 \(z=e^{u+i\theta}\)（\(u>0\)，\(\theta\in(-\pi,\pi]\)）。
由 \(w=z+z^{-1}\) 得
\[
w=2\cosh u\cos\theta+2i\sinh u\sin\theta,
\]
即坐标表达式。
\par
固定 \(u\) 时，
\[
\left(\frac{x}{2\cosh u}\right)^2+\left(\frac{y}{2\sinh u}\right)^2=1,
\]
其焦距为 \(2\sqrt{(2\cosh u)^2-(2\sinh u)^2}=4\)，故焦点为 \(\pm2\)。
固定 \(\theta\) 且 \(\sin\theta\cos\theta\neq 0\) 时，
\[
\left(\frac{x}{2\cos\theta}\right)^2-\left(\frac{y}{2\sin\theta}\right)^2=1,
\qquad
(2\cos\theta)^2+(2\sin\theta)^2=4,
\]
同样以 \(\pm2\) 为焦点。
\par
计算导数
\[
(x_u,y_u)=(2\sinh u\cos\theta,\ 2\cosh u\sin\theta),\qquad
(x_\theta,y_\theta)=(-2\cosh u\sin\theta,\ 2\sinh u\cos\theta),
\]
得内积 \(x_u x_\theta+y_u y_\theta=0\)，并且
\[
x_u^2+y_u^2=x_\theta^2+y_\theta^2=4(\cosh^2u-\cos^2\theta),
\]
从而得到度量表达式。证毕。
\end{proof}

\leanverified{paper\_conclusion\_jg\_focal\_distance\_decoding}
\begin{theorem}[焦点差分--加和的相位--尺度解码]\label{thm:conclusion-jg-focal-distance-decoding}
对任意 \(w\in\Omega_{\mathrm{slit}}\) 定义两焦距
\[
d_{+}(w):=|w-2|,
\qquad
d_{-}(w):=|w+2|.
\]
令 \((u,\theta)\) 为定理 \ref{thm:conclusion-jg-slit-uniformization-orthogonal-coordinates} 给出的参数（取 \(\theta\in[0,\pi]\) 的代表）。
则有显式解码公式
\[
\cosh u=\frac{d_{+}(w)+d_{-}(w)}{4},
\qquad
\cos\theta=\frac{d_{-}(w)-d_{+}(w)}{4}.
\]
进一步令 \(r:=e^{u}\ge 1\)，则 \(r\) 满足
\[
r+\frac1r=\frac{d_{+}(w)+d_{-}(w)}{2},
\]
从而尺度 \(u=\log r\) 与相位 \(\theta\) 均可由单点 \(w\) 唯一确定。
\end{theorem}
\begin{proof}
令 \(z=z(w)=e^{u+i\theta}\) 为定理 \ref{thm:conclusion-jg-slit-uniformization-orthogonal-coordinates} 的均匀化根（\(|z|>1\)）。
由恒等式
\[
w-2=z+z^{-1}-2=\frac{(z-1)^2}{z},\qquad
w+2=z+z^{-1}+2=\frac{(z+1)^2}{z},
\]
得
\[
d_{+}(w)=\frac{|z-1|^2}{|z|},\qquad
d_{-}(w)=\frac{|z+1|^2}{|z|}.
\]
写 \(z=re^{i\theta}\)（\(r>1\)），则
\[
|z\pm1|^2=r^2+1\pm2r\cos\theta.
\]
故
\[
d_{-}(w)-d_{+}(w)=\frac{4r\cos\theta}{r}=4\cos\theta,
\qquad
d_{-}(w)+d_{+}(w)=\frac{2(r^2+1)}{r}=2\left(r+\frac1r\right)=4\cosh(\log r),
\]
即得到所示解码公式与 \(r+r^{-1}\) 的恒等式。证毕。
\end{proof}

\begin{proposition}[两素数 Gödel 尺度的对数格化与逼近度量]\label{prop:conclusion-jg-godel-loglattice-approximation}
取互异素数 \(p\neq q\)，并定义乘法群
\[
H:=\{p^{m}q^{n}:\ (m,n)\in\ZZ^2\}\subset\RR_{>0},
\qquad
\Lambda:=\log H=\{m\log p+n\log q:\ (m,n)\in\ZZ^2\}\subset\RR.
\]
对任意 \(w\in\Omega_{\mathrm{slit}}\)，令 \(u(w):=\log|z(w)|>0\) 为定理 \ref{thm:conclusion-jg-slit-uniformization-orthogonal-coordinates} 的尺度参数。
则：
\begin{enumerate}
\normalfont
  \item \(e^{u(w)}\in H\) 当且仅当 \(u(w)\in\Lambda\)；
  \item 并且
  \[
  \inf_{h\in H}\left|\log\frac{e^{u(w)}}{h}\right|
  =\operatorname{dist}\bigl(u(w),\Lambda\bigr),
  \qquad
  \inf_{h\in H}\frac{e^{u(w)}}{h}
  =\exp\!\bigl(\operatorname{dist}(u(w),\Lambda)\bigr).
  \]
\end{enumerate}
\end{proposition}
\begin{proof}
第一条由定义即得。
第二条令 \(u=u(w)\)，并记 \(h=e^\ell\)（\(\ell\in\Lambda\)），则
\(
|\log(e^u/h)|=|u-\ell|
\)；
对 \(\ell\in\Lambda\) 取下确界得到距离公式，指数化得第二式。证毕。
\end{proof}
\leanverified{paper\_conclusion\_jg\_godel\_loglattice\_approximation}

\begin{theorem}[两素数 Gödel 尺度在双曲叶上的稠密轨道]\label{thm:conclusion-jg-godel-dense-orbits-on-hyperbola-leaf}
\leanverified{paper\_conclusion\_jg\_godel\_dense\_orbits\_on\_hyperbola\_leaf}
取互异素数 \(p\neq q\)，并令 \(H\) 如命题 \ref{prop:conclusion-jg-godel-loglattice-approximation}。
固定 \(\theta\in(0,\pi)\) 且 \(\sin\theta\cos\theta\neq 0\)，定义双曲叶
\[
\mathcal H_{\theta}:=\{\,w(r):=r e^{i\theta}+r^{-1}e^{-i\theta}:\ r>0\,\}\subset\Omega_{\mathrm{slit}}.
\]
则集合
\[
\mathcal O_{\theta}:=\{\,w(r):\ r\in H\,\}
\]
在 \(\mathcal H_{\theta}\) 中稠密（相对拓扑）。
\end{theorem}
\begin{proof}
若 \(\log p/\log q\in\QQ\)，则存在 \(m,n\neq 0\) 使 \(p^m=q^n\)，与唯一分解矛盾，故 \(\log p/\log q\notin\QQ\)。
于是加法群 \(\Lambda=\{m\log p+n\log q\}\) 在 \(\RR\) 中稠密，从而 \(H=\exp(\Lambda)\) 在 \(\RR_{>0}\) 中稠密。
映射 \(r\mapsto w(r)\) 连续，故 \(\mathcal O_\theta\) 在 \(\mathcal H_\theta\) 中稠密。证毕。
\end{proof}

\begin{theorem}[有理余弦网格在区间 \(0\le \theta \le \pi\) 上的稠密性]\label{thm:conclusion-hypercube-phase-grid-dense}
\leanverified{paper\_conclusion\_hypercube\_phase\_grid\_dense}
令
\[
S:=\left\{\theta_{d,k}:=\arccos\!\left(1-\frac{2k}{d}\right):\ d\in\NN,\ 0\le k\le d\right\}\subset[0,\pi].
\]
则 \(S\) 在 \([0,\pi]\) 中稠密。
\end{theorem}
\begin{proof}
集合 \(\{1-2k/d\}\) 为 \([-1,1]\) 中的有理点族，随 \(d\) 变化在 \([-1,1]\) 中稠密。
函数 \(\arccos:[-1,1]\to[0,\pi]\) 连续且为同胚，故稠密性在像下保持。证毕。
\end{proof}

\begin{theorem}[Gödel 尺度与相位网格耦合的裂平面稠密采样]\label{thm:conclusion-jg-godel-hypercube-dense-sampling-omega}
\leanverified{paper\_conclusion\_jg\_godel\_hypercube\_dense\_sampling\_omega}
取互异素数 \(p\neq q\)，并令 \(H\) 如命题 \ref{prop:conclusion-jg-godel-loglattice-approximation}，令 \(S\) 如定理 \ref{thm:conclusion-hypercube-phase-grid-dense}。
定义采样集
\[
\mathcal D:=\{\,r e^{i\theta}+r^{-1}e^{-i\theta}:\ r\in H,\ \theta\in S\,\}\subset\Omega_{\mathrm{slit}}.
\]
则 \(\mathcal D\) 在 \(\Omega_{\mathrm{slit}}\) 中稠密。
\end{theorem}
\begin{proof}
由定理 \ref{thm:conclusion-jg-slit-uniformization-orthogonal-coordinates}，任意 \(w\in\Omega_{\mathrm{slit}}\) 都可写作
\(
w=r e^{i\theta}+r^{-1}e^{-i\theta}
\)
其中 \(r=e^{u}>1\)、\(\theta\in(0,\pi)\)。
由 \(H\) 在 \(\RR_{>0}\) 中稠密可取 \(r_n\in H\) 使 \(r_n\to r\)；
由定理 \ref{thm:conclusion-hypercube-phase-grid-dense} 可取 \(\theta_n\in S\) 使 \(\theta_n\to\theta\)。
连续性给出
\(
r_n e^{i\theta_n}+r_n^{-1}e^{-i\theta_n}\to w
\)，故 \(\mathcal D\) 稠密。证毕。
\end{proof}

\begin{theorem}[Lee--Yang 单位圆零点在 Gödel 尺度变参下的双曲闭包]\label{thm:conclusion-leyang-godel-hyperbola-orbit-closure}
设
\[
P(z)=\prod_{j=1}^{N}(z-e^{i\theta_j}),\qquad \theta_j\in[0,2\pi),
\]
且其全部零点位于单位圆 \(|z|=1\)。
对 \(r>0\) 定义点层面的输运
\[
w_j(r):=r e^{i\theta_j}+r^{-1}e^{-i\theta_j}.
\]
取互异素数 \(p\neq q\)，令 \(H=\{p^{m}q^{n}\}\)。
则对每个满足 \(\sin\theta_j\cos\theta_j\neq 0\) 的 \(j\)，集合
\[
\{\,w_j(r):\ r\in H\,\}
\]
在双曲叶
\(
\mathcal H_{\theta_j}=\{w_j(r):r>0\}
\)
中稠密；并且在 \(\Omega_{\mathrm{slit}}\) 的相对拓扑下，
\[
\overline{\bigcup_{j=1}^{N}\{\,w_j(r):\ r\in H\,\}}
=\bigcup_{j=1}^{N}\mathcal H_{\theta_j}.
\]
\end{theorem}
\begin{proof}
对固定 \(j\)，映射 \(r\mapsto w_j(r)\) 连续且其像为 \(\mathcal H_{\theta_j}\)。
由定理 \ref{thm:conclusion-jg-godel-dense-orbits-on-hyperbola-leaf} 得 \(\{w_j(r):r\in H\}\) 在 \(\mathcal H_{\theta_j}\) 中稠密。
有限并的闭包为闭包之并，故得到整体闭包恒等式。证毕。
\end{proof}

\begin{conjecture}[对数素数格的最短向量控制相干阈值]\label{conj:conclusion-jg-godel-coherence-threshold-short-vector}
取互异素数 \(p\neq q\)，并令 \(\Lambda=\{m\log p+n\log q\}\subset\RR\)。
对任意 \(w\in\Omega_{\mathrm{slit}}\)，令 \(u(w)=\log|z(w)|\) 为尺度读出，并定义相干缺陷
\[
\Delta(w):=\operatorname{dist}\bigl(u(w),\Lambda\bigr).
\]
预期存在一类由固定双曲叶 \(\mathcal H_\theta\) 与尺度坐标 \(u\) 所生成的自然判别/识别问题，使得其临界常数仅依赖于 \(\Lambda\) 的几何数论不变量（等价地，依赖于 \(\log p/\log q\) 的最佳丢番图逼近常数），并且在定理 \ref{thm:conclusion-leyang-godel-hyperbola-orbit-closure} 的 Lee--Yang 输运闭包与定理 \ref{thm:conclusion-hypercube-phase-grid-dense} 的相位离散化两种机制下保持一致。
\end{conjecture}

\begin{theorem}[Joukowsky--Godel 输运的二次提升结构与 wreath 嵌入]\label{thm:conclusion-joukowsky-transport-wreath-embedding}
设 \(K\) 为特征不等于 \(2\) 的域，\(r\in K^\times\)。
取次数 \(n\) 的可分多项式 \(P(z)\in K[z]\)，并假设 \(P(0)\neq 0\)。
定义输运多项式
\[
Q_r(w):=\operatorname{Res}_z\!\bigl(P(z),\,r z^2-wz+r^{-1}\bigr)\in K[w].
\]
记 \(w_1,\dots,w_m\) 为 \(Q_r\) 的互异根，\(M\) 为 \(Q_r\) 在 \(K(r)\) 上的分裂域，\(L\) 为 \(P\) 在 \(K(r)\) 上的分裂域。
则有域包含
\[
L\subset
M\!\left(\sqrt{w_1^2-4},\dots,\sqrt{w_m^2-4}\right),
\]
因而
\[
[L:M]\mid 2^m\le 2^n.
\]
并且自然诱导群嵌入
\[
\mathrm{Gal}(L/K(r))
\hookrightarrow
(\ZZ/2\ZZ)^m\wr \mathrm{Gal}(M/K(r)).
\]
\end{theorem}
\begin{proof}
若 \(z\) 为 \(P\) 的根，则
\[
w:=J_r(z)=rz+r^{-1}z^{-1}
\]
使得 \(Q_r(w)=0\)。
固定 \(Q_r\) 的任一根 \(w_i\)，对应 \(z\) 满足
\[
r z^2-w_i z+r^{-1}=0,
\]
其判别式为 \(\Delta_i=w_i^2-4\)，故
\[
z=\frac{w_i\pm\sqrt{w_i^2-4}}{2r}.
\]
因此每个 \(P\) 的根都落在某个
\(
M(\sqrt{w_i^2-4})
\)
中，进而
\[
L\subset M\!\left(\sqrt{w_1^2-4},\dots,\sqrt{w_m^2-4}\right).
\]
故扩张次数整除 \(2^m\)，并由 \(m\le n\) 得 \([L:M]\mid 2^m\le 2^n\)。

记
\[
N:=M\!\left(\sqrt{w_1^2-4},\dots,\sqrt{w_m^2-4}\right).
\]
在 \(N/M\) 上，每个 \(\sqrt{w_i^2-4}\mapsto -\sqrt{w_i^2-4}\) 生成一个对合，故
\[
\mathrm{Gal}(N/M)\le (\ZZ/2\ZZ)^m.
\]
同时 \(\mathrm{Gal}(M/K(r))\) 在根集 \(\{w_i\}\) 上的置换作用会诱导对上述 \(m\) 个对合坐标的置换作用，从而得到半直积结构
\[
\mathrm{Gal}(N/K(r))
\hookrightarrow
(\ZZ/2\ZZ)^m\rtimes \mathrm{Gal}(M/K(r))
\hookrightarrow
(\ZZ/2\ZZ)^m\wr \mathrm{Gal}(M/K(r)).
\]
由于 \(L\subset N\)，限制映射给出
\[
\mathrm{Gal}(L/K(r))\hookrightarrow \mathrm{Gal}(N/K(r)),
\]
合并即得结论。证毕。
\end{proof}

\begin{corollary}[输运全息的最小离散账本下界]\label{cor:conclusion-joukowsky-discrete-ledger-lower-bound}
\leanverified{paper\_conclusion\_joukowsky\_discrete\_ledger\_lower\_bound}
沿用定理 \ref{thm:conclusion-joukowsky-transport-wreath-embedding} 的记号，令
\[
H:=\mathrm{Gal}(L/M)\le (\ZZ/2\ZZ)^m.
\]
在固定 \((Q_r,r)\) 的条件下，与之相容的 \(P\) 至少分裂为 \(|H|\) 个互不等价的代数分支。
因此任何从 \((Q_r,r)\) 确定性恢复 \(P\) 的方案都必须附加离散账本，且其信息量下界为
\[
\log_2|H|.
\]
若判别式类 \(\{w_i^2-4\}\) 在 \(M^\times/(M^\times)^2\) 中独立（\(\FF_2\)-秩为 \(m\)），则 \(|H|=2^m\)，最小账本达到 \(m\) 比特。
\end{corollary}
\begin{proof}
\(H\) 的元素固定 \(M\)，因而固定全部 \(w_i\)，但可独立翻转若干 \(\sqrt{w_i^2-4}\) 的符号。
在
\[
z=\frac{w_i\pm\sqrt{w_i^2-4}}{2r}
\]
中，这等价于交换同一 \(w_i\) 下的二次提升根对。
因此同一 \((Q_r,r)\) 的前像至少包含 \(|H|\) 个不可区分的代数分支。
若恢复算法不额外输入可区分这 \(|H|\) 个分支的离散信息，则无法确定唯一前像，所需最小信息量即 \(\log_2|H|\)。
当判别式类独立时，翻转自由度满秩，故 \(|H|=2^m\)。证毕。
\end{proof}

\begin{theorem}[6D 折叠压缩比的椭圆化与黄金比指数实现]\label{thm:conclusion-6d-folding-golden-ellipse-realization}
记
\[
I_m:=\frac{|\Omega_m|}{|X_m|}=\frac{2^m}{F_{m+2}},
\qquad
r_m:=\sqrt{I_m},
\]
其中 \((F_m)\) 为 Fibonacci 数列，\(\varphi=(1+\sqrt5)/2\)。
则有：
\begin{enumerate}
  \item 指数斜率满足
  \[
  \frac{1}{m}\log r_m^2\longrightarrow \log\!\left(\frac{2}{\varphi}\right)\qquad(m\to\infty).
  \]
  \item 若 \(\mu_{r_m}:=(J_{r_m})_\#m\) 支撑于椭圆 \(E_{r_m}\)，则
  \[
  \int_{E_{r_m}}|w|^2\,\dd\mu_{r_m}(w)
  =r_m^2+r_m^{-2}
  =I_m+I_m^{-1},
  \]
  从而
  \[
  \log\!\int_{E_{r_m}}|w|^2\,\dd\mu_{r_m}(w)=\log I_m+o(1),
  \]
  主阶斜率同为 \(\log(2/\varphi)\)。
  \item 在 \(m=6\) 时，
  \[
  r_6=\sqrt{\frac{64}{21}},
  \qquad
  \int_{E_{r_6}}|w|^2\,\dd\mu_{r_6}(w)
  =\frac{64}{21}+\frac{21}{64}.
  \]
\end{enumerate}
\end{theorem}
\begin{proof}
由 Binet 公式 \(F_{m+2}=\varphi^{m+2}/\sqrt5+O(\varphi^{-m})\) 得
\[
\log I_m
=m\log\!\left(\frac{2}{\varphi}\right)+O(1),
\]
即
\[
\frac1m\log r_m^2=\frac1m\log I_m\to \log\!\left(\frac{2}{\varphi}\right).
\]
第二部分由
\(
|r_m e^{i\theta}+r_m^{-1}e^{-i\theta}|^2
=r_m^2+r_m^{-2}+e^{2i\theta}+e^{-2i\theta}
\)
对 \(\theta\) 平均得到。
且 \(I_m\to\infty\)，故
\[
\log(I_m+I_m^{-1})=\log I_m+o(1).
\]
第三部分代入 \(m=6\) 与 \(F_8=21\) 即得。证毕。
\end{proof}

\begin{theorem}[时间视窗的圆维下界与残差账本增长律]\label{thm:conclusion-time-window-cdim-ledger-lower-bound}
设时间视窗群 \(G_T:=\ZZ^T\)，并取 \(b\in\NN\)。
定义 \(b\)-位窗口对象
\[
W_b(G_T):=(\ZZ/2^b\ZZ)^T,\qquad |W_b(G_T)|=2^{bT}.
\]
设存在可逆编码
\[
E_{T,b}:W_b(G_T)\hookrightarrow \mathrm{Phase}(k)_b\times R_{T,b},
\]
其中 \(|\mathrm{Phase}(k)_b|=2^{bk}\)，\(R_{T,b}\) 为残差账本。
则必有
\[
|R_{T,b}|\ge 2^{(T-k)b}.
\]
特别地：
\begin{enumerate}
  \item 若 \(k\) 不随 \(T\) 增长，则 \(|R_{T,b}|\) 随 \(T\) 指数增长；
  \item 若 \(b\) 固定且要求 \(|R_{T,b}|\le T^C\)（某常数 \(C>0\)），则
  \[
  k\ge T-\frac{C}{b}\log_2 T;
  \]
  \item 若进一步 \(b\gtrsim \log T\)，则上式强化为 \(k\ge T-O(1)\)。
\end{enumerate}
\end{theorem}
\begin{proof}
由单射 \(E_{T,b}\) 得计数不等式
\[
2^{bT}=|W_b(G_T)|
\le
|\mathrm{Phase}(k)_b|\,|R_{T,b}|
=
2^{bk}|R_{T,b}|,
\]
故
\[
|R_{T,b}|\ge 2^{(T-k)b}.
\]
第 (1) 条直接成立。
若 \(|R_{T,b}|\le T^C\)，则
\[
2^{(T-k)b}\le T^C
\quad\Longrightarrow\quad
T-k\le \frac{C}{b}\log_2 T,
\]
得到第 (2) 条。
当 \(b\gtrsim \log T\) 时，右端有界，故 \(k\ge T-O(1)\)。证毕。
\end{proof}
\leanverified{paper_conclusion_time_window_cdim_ledger_lower_bound}

\begin{remark}[解析全息、代数提升、压缩几何与时间圆维的闭环]\label{rem:conclusion-holography-uplift-ellipse-time-closed-loop}
定理 \ref{thm:conclusion-analytic-holography-inversion-rigidity} 给出解析全息条件下的几何反演刚性与不可见参数的唯一非解析读出；
定理 \ref{thm:conclusion-joukowsky-transport-wreath-embedding} 与推论 \ref{cor:conclusion-joukowsky-discrete-ledger-lower-bound}
把输运前像的不唯一性精确压缩为二次提升核的离散账本维度；
定理 \ref{thm:conclusion-6d-folding-golden-ellipse-realization}
将 \(6\) 维窗口压缩率与黄金比指数编码为椭圆二阶径向矩；
定理 \ref{thm:conclusion-time-window-cdim-ledger-lower-bound}
则在时间视窗方向给出固定相位维压缩所必需承担的残差下界。
四者共同构成“解析全息 \(\to\) 代数分歧 \(\to\) 几何压缩 \(\to\) 时间圆维约束”的闭合审计链条。
\end{remark}
